\documentclass[11pt]{article}
\usepackage[margin=1in]{geometry}
\usepackage[utf8]{inputenc}
\usepackage[T1]{fontenc}
\usepackage{enumitem}
\usepackage{hyperref}
\usepackage{xcolor}
\usepackage{titlesec}

\titlespacing*{\section}{0pt}{8pt}{4pt}
\titlespacing*{\subsection}{0pt}{6pt}{3pt}
\setlength{\parskip}{4pt}
\setlength{\parindent}{0pt}
\pagestyle{empty}

\begin{document}

\begin{center}
{\Large\bfseries Collaboration Brief: Google.org AI for Science Application}\\[4pt]
{\normalsize Patrick Cooper, Department of Computer Science, CU Boulder}\\[2pt]
{\small April 2026}
\end{center}

\vspace{-4pt}
\hrule
\vspace{6pt}

\section*{The Opportunity}

The \href{https://www.google.org/impact-challenges/ai-science/}{Google.org Impact Challenge: AI for Science} is a \$30M global open call offering grants of \$500K--\$3M to academic institutions. Selected teams also receive six months of pro bono Google engineering support and Google Cloud credits. \textbf{Applications close April 17, 2026.} The two priority focus areas are \emph{AI for Health \& Life Sciences} and \emph{AI for Climate Resilience \& Environmental Science}.

\section*{What We Have Built}

Over the past year at CU Boulder, I have developed \textsc{DeFAb} (Defeasible Abduction Benchmark): an open-source pipeline and dataset for evaluating and training foundation models on \emph{exception-aware reasoning} over structured knowledge bases. The core idea is simple: medical and scientific knowledge is full of rules that hold \emph{by default} but admit exceptions. Drug X treats Condition Y, unless the patient has Comorbidity Z. Gene G has molecular function F, unless NOT-qualified experimental evidence overrides that annotation. SNOMED concept C inherits properties from its parent, unless the child explicitly overrides them.

Current AI benchmarks (MedQA, PubMedQA, BioASQ) test whether models can \emph{retrieve} biomedical facts. None test whether models can \emph{reason about exceptions}: recognizing when a general rule fails, identifying what exception applies, and revising conclusions while preserving unrelated knowledge. This is the gap between an AI that looks up answers and one that reasons safely.

\textsc{DeFAb} addresses this gap with a formally grounded methodology. We convert expert-curated knowledge bases into \emph{defeasible theories} where every rule is tagged as strict or revisable, every conclusion has a traceable proof, and every exception is verifiable in polynomial time. We then generate evaluation instances at three levels of difficulty: (1)~missing-fact completion, (2)~rule abduction, and (3)~\emph{defeater abduction}, where the model must construct an exception that blocks an incorrect default conclusion while preserving all unrelated inferences.

\textbf{Critically, we have already built extraction pipelines for the major biomedical ontologies:}
\begin{itemize}[nosep, leftmargin=1.5em]
    \item \textbf{Gene Ontology:} 409,378 rules extracted; 1,250 defeaters from NOT-qualified experimental annotations
    \item \textbf{UMLS:} Extractor implemented for 3.4M concepts across 189 source vocabularies; encodes clinical contradiction relations (e.g., \texttt{treats} vs.\ \texttt{contraindicated\_with}) as structured defeaters
    \item \textbf{SNOMED CT:} Extractor implemented for 360K+ clinical concepts; captures property-override exceptions via OWL axioms
    \item \textbf{MeSH:} 76,650 strict taxonomic rules across 16 medical subject categories
\end{itemize}

Preliminary evaluation of four frontier models (GPT-5.2, Claude Sonnet 4.6, DeepSeek-R1, Kimi-K2.5) shows strong performance on rule retrieval but significant failures on exception handling, precisely the capability that matters for safe biomedical inference.

\section*{The Proposed Application: DeFAb-Health}

We propose to apply under \emph{AI for Health \& Life Sciences} with a project titled \textbf{DeFAb-Health: A Verifier-Backed Benchmark for Exception-Aware Biomedical AI}. The deliverables:

\begin{enumerate}[nosep, leftmargin=1.5em]
    \item An \textbf{open benchmark} of defeasible reasoning instances drawn from GO, UMLS, SNOMED~CT, and MeSH, testing whether AI systems can handle the exceptions inherent in biomedical knowledge
    \item A \textbf{verifier-backed training pipeline} that uses polynomial-time proof checking as an exact reward signal, enabling reliable fine-tuning without human annotation
    \item A \textbf{public leaderboard} tracking foundation model performance on biomedical exception handling
    \item Demonstrated \textbf{fine-tuned models} with measurably improved defeasible reasoning on biomedical theories
\end{enumerate}

The grant's criteria require \emph{innovative AI use}, \emph{open-source release}, \emph{feasibility} (existing traction matters), and \emph{scalability}. We satisfy all four: AI is the core methodology, everything will be open-source, the pipeline is already built and tested, and the extraction architecture is ontology-agnostic, so any institution can extend it to their domain.

\section*{Why Collaboration with Anschutz Is Essential}

The grant strongly favors interdisciplinary teams with domain expertise. A CS-only proposal lacks credibility on \emph{why these reasoning failures matter clinically}. We need collaborators who can:

\begin{itemize}[nosep, leftmargin=1.5em]
    \item Validate that exception-handling failures in UMLS/SNOMED reasoning represent real clinical risk
    \item Help design evaluation scenarios grounded in actual biomedical inference tasks (drug interaction checking, clinical decision support, genomic annotation interpretation)
    \item Provide or facilitate access to licensed biomedical data (UMLS, SNOMED CT)
    \item Co-author the application with domain authority
\end{itemize}

The time constraint is tight: two weeks to submission. The core infrastructure is built. What we need is the biomedical grounding to make the health impact case convincing to reviewers.

\vspace{4pt}
\hrule
\vspace{4pt}
{\small\textbf{Contact:} Patrick Cooper, \href{mailto:patrick.cooper@colorado.edu}{patrick.cooper@colorado.edu} \hfill \textbf{Project:} \href{https://github.com/PatrickAllenCooper/blanc}{github.com/PatrickAllenCooper/blanc} (open-source, MIT license)}

\end{document}
